\subsection{引言}\label{sec:3.1}
统计分布用来为总体建立模型；因此，我们通常处理的是一个分布族（family of distributions），而不是单个分布. 这个族由一个或多个参数标定，使我们可以在保持函数形式不变的同时，改变分布的某些特征. 例如，我们可能认为正态分布适合为某个总体建模，但不能准确指定其均值. 此时处理的是一个参数族：均值为 $\mu$ 的正态分布，其中 $\mu$ 是尚未指定的参数，$-\infty<\mu<\infty$.

本章罗列许多较常见的统计分布，其中一些已经遇到过. 对每个分布，我们将给出均值、方差以及许多有助于理解的其他有用或描述性的度量. 我们还会介绍这些分布的一些典型应用，以及它们之间某些有趣而有用的联系. 部分事实汇总于书末的表格中. 本章对统计分布的介绍远非全面；这一任务已由 \hyperref[bib:source177]{Johnson and Kotz (1969–1972)} 的多卷本著作 \textit{Distributions in Statistics}，以及更新版 Johnson, Kotz, and Balakrishnan (\hyperref[bib:source178]{1994}, \hyperref[bib:source179]{1995}) 和 \hyperref[bib:source180]{Johnson, Kotz, and Kemp (1992)} 完成.
